% บทคัดย่อ -- เขียนหลังทำบทที่ 4 และ 5 เสร็จ ความยาวไม่เกิน 1 หน้า
% สูตรย่อ: ทำอะไร -> เก็บข้อมูลอย่างไร -> พบอะไร (ใส่ตัวเลขจริง) -> แปลว่าอะไร
\thispagestyle{empty}

\begin{center}
    {\fontsize{16pt}{20pt}\selectfont\bfseries บทคัดย่อ\par}
\end{center}

\vspace{1em}

% TODO: ย่อหน้าที่ 1 -- ที่มาและวัตถุประสงค์ 3--4 บรรทัด

% TODO: ย่อหน้าที่ 2 -- วิธีดำเนินการ: กลุ่มตัวอย่าง 121 คน แบบสอบถามออนไลน์
% DASS-21 ฉบับแปลไทย สถิติที่ใช้ (t-test, ANOVA, สหสัมพันธ์, สหสัมพันธ์บางส่วน)

% TODO: ย่อหน้าที่ 3 -- ผลที่พบ พร้อมตัวเลข เช่น
% เวลาใช้โซเชียลมีเดียไม่สัมพันธ์กับมิติใดของ DASS-21 อย่างมีนัยสำคัญ
% การเสพข่าวสัมพันธ์กับความวิตกกังวล และสภาพการเงินสัมพันธ์กับภาวะซึมเศร้ามากที่สุด

% TODO: ย่อหน้าที่ 4 -- ข้อสรุปและข้อเสนอแนะสั้น ๆ

\vspace{1em}

\noindent\textbf{คำสำคัญ:} โซเชียลมีเดีย, สุขภาวะทางจิต, DASS-21, นักศึกษา, ความเพียงพอทางการเงิน

\clearpage
